\chapter{Conclusiones y estado de cierre}

La solución acumulativa conserva la VPC y los controles demostrados en la Fase 1 y añade una DMZ,
DNS interno, intranet de Recursos Humanos, web pública y clientes temporales de Visitas y TI. La
infraestructura permanece reproducible en Terraform y el cómputo de la Fase 2 está deshabilitado de
forma predeterminada para prevenir cargos accidentales.

La revisión local comprobó formato y validez de Terraform, sintaxis de plantillas, reglas de
retorno de las NACL, límites de los Security Groups y 50 flujos de la matriz estática. Los clientes
temporales instalan \code{nftables}; TI usa el Gateway Endpoint de S3 para obtener el paquete sin
salida general a Internet y valida la resolución interna antes de declararse listo.

Andrés Rafael Chivalán ejecutó la ventana AWS y recopiló las evidencias; Pablo Daniel Barillas
Moreno revisó e integró el material. La suite registró 36 subpruebas correctas y 3 omitidas, sin
contradicciones. Las pruebas aplicables demostraron resolución DNS interna, intranet privada,
navegación de Visitas,
acceso iniciado por TI, bloqueo de conexiones hacia TI y exposición pública limitada de la DMZ.

La auditoría posterior a la ventana registró cero instancias EC2 y cero recursos etiquetados con el
proyecto en \code{us-east-1}. La VPC \code{172.31.0.0/16} que aparece en el inventario es la VPC
predeterminada de la cuenta y no pertenece a RedNova.

\begin{figure}[H]
  \centering
  \includegraphics[width=0.92\textwidth]{04-inventario-aws-posterior-al-cierre.png}
  \caption{Inventario posterior al cierre: sin Elastic IP, sin recursos con la etiqueta del
  proyecto y sin instancias EC2 en las regiones consultadas.}
  \label{fig:f2-cierre-aws}
\end{figure}

\begin{warningbox}[Pendientes de entrega, no de implementación base]
Antes de entregar todavía debe confirmarse la aceptación docente de Security Groups más NACL
como firewall y ACL del ejercicio. La captura dedicada de la página
pública por Elastic IP ya está incorporada. HTTPS y la redirección permanecen como punto extra
opcional.
\end{warningbox}
